\subsubsection{为何本卷必读本章}

算法理论只是冰山一角。赛场成败往往取决于工程习惯与协作方式。
下列五条凝练自队内调车经验（源自自瞄核心算法讲义结语），开营与结营各读一次；培训负责人可据此做过程检查。

\begin{example}[把“好了”改写成证据]
原记录是“今天调好了”。改成“固定录包 A，配置版本 B，目标帧范围 100--150，修改前位置有三次跳变；本次只调整角点筛选，修改后同区间无跳变，尚未实车验证”。检查：他人可以找到输入、修改和比较对象，且不会误把回放成功当作已通过实车。
\end{example}

\begin{exercise}
\label{v3-ex-engineering-1}
一项改动只在仿真中通过，记录中应怎样写结论与下一步？
\end{exercise}

\subsubsection{五条准则}

\begin{enumerate}
    \item \textbf{实测优先，仿真为辅。}
    仿真难以复现真实光照、振动、弹道散布与对手机动。新算法须经靶场与对抗演练再定稿。

    \item \textbf{边测边写，过程正确比结果正确更重要。}
    不要只盯最终命中率——偶然命中可能是「蒙对」。应逐环节检查：检测是否稳定、解算是否跳变、滤波是否收敛、控制是否平滑。每一步可复现，整体才可复现。

    \item \textbf{拒绝重复造轮子；重构是学习，调车是工作。}
    成熟的 PnP、Kalman、MPC、检测框架优先复用。为理解原理而重构关键模块属于\textbf{学习时间}；临场调车属于\textbf{工程时间}：稳定优先，用验证过的组件。

    \item \textbf{强化四项意识：时间、进度、合作、钻研。}
    视觉跨识别、估计、控制、机械接口，单人难通吃。设里程碑、定期同步消息与坐标系；难点查资料或请教，避免无效死磕。

    \item \textbf{先定位问题，再寻找解，切忌盲目乱试。}
    偏弹、抖动、跳变时，先用日志、可视化与离线回放锁定异常环节与时刻，再改参数或算法。随机拧参「蒙好」会在下一场暴露更多隐患。
\end{enumerate}

\begin{example}[按五条准则安排一次修改]
教学故障是扫场时位置跳变。先在现成录包上复现，逐段查看原图与角点；保留已验证的解算器，仅修改角点筛选；分配采集、分析与复核负责人；写下假设后做同一录包的对照。检查：五条准则分别落在证据、过程、复用、协作和定位上。
\end{example}

\begin{exercise}
\label{v3-ex-engineering-2}
临场前一天想重写已验证的滤波器，理由是“自己写更懂”。应如何安排学习与工程工作？
\end{exercise}

\subsubsection{过程检查表示意}

\begin{center}
\begin{tabular}{@{}p{0.72\textwidth}p{0.20\textwidth}@{}}
\toprule[1.5pt]
检查项 & 是 / 否 \\
\midrule[1pt]
本周改动是否有实车或录包证据 &  \\
是否能指出问题发生在检测/解算/滤波/控制哪一段 &  \\
接口变更是否更新了文档与消息定义 &  \\
是否区分了学习实验分支与比赛主分支 &  \\
\bottomrule[1.5pt]
\end{tabular}
\end{center}

\begin{example}[填写一条失败记录]
按检查表填入：输入录包 A、基线版本 B、单一改动 C、比较窗口和指标、结果恶化、恢复 B。即使改动失败，记录仍能排除一个假设。检查：失败结果保留，接口变更同时记录发送端和接收端验证。
\end{example}

\begin{note}
实车运动与开火测试须在队内规定场地和监督条件下进行，先确认急停、限位及人员位置。离线回放通过只能支持其覆盖的条件，不能直接记作实车已通过。
\end{note}

\begin{conclusion}
\begin{itemize}
  \item 固定基线，写下假设。
  \item 一次改变一个主要因素，比较同一输入下的中间量与结果。
  \item 保存失败实验与回滚记录。
\end{itemize}
\end{conclusion}

\begin{exercise}
\label{v3-ex-engineering-3}
测试结果变差后，删掉失败记录只留“最终好参数”有什么问题？
\end{exercise}

\subsubsection{适用范围}

五条不仅适用于自瞄，也适用于雷达站多进程联调与导航上图。从「能跑」到「能战」，靠的是可复现的过程，而不是单次好运。

\begin{example}[把准则迁移到导航]
把自瞄记录中的原图、角点和控制量替换为地图、位姿、路径和速度，仍按同样顺序记录版本、场景、最早异常与单一改动。检查：准则迁移的是取证方法，不能把自瞄阈值照搬到导航。
\end{example}

\begin{exercise}
\label{v3-ex-engineering-4}
雷达通信偶发积压时，“先定位再寻找解”应怎样落实？给出两个可观察量。
\end{exercise}
